\documentclass[11pt]{article}
\usepackage{amsmath,amssymb,hyperref}
\usepackage{geometry}
\usepackage{url}
\usepackage{booktabs}
\Urlmuskip=0mu plus 1mu
\geometry{margin=1in}

% Write-up for Experiments/SU4_SoV_Orthogonality_Untwisted_L2.wb.
% Twisted construction: ../Clean/SU4_SoV_ScalarProduct_L2.tex (and .wb).

\newcommand{\Op}{\mathcal{O}}
\newcommand{\shift}[1]{^{[#1]}}

\title{Orthogonality of separated-variables states at zero twist\\
for the $\mathfrak{gl}(4)$ spin chain with rectangular representations}
\author{Notes accompanying \texttt{Experiments/SU4\_SoV\_Orthogonality\_Untwisted\_L2.wb}}
\date{11 September 2026}

\begin{document}
\maketitle

\begin{abstract}
For the twisted $\mathfrak{gl}(4)$ chain of length $L=2$ with sites $[S,S,0,0]\otimes[P,P,0,0]$,
orthogonality of transfer-matrix eigenstates follows from a linear system in eight differences of
integrals of motion. We repeat the argument at zero twist. There, six of the eight differences
become linear combinations of the remaining two, which are the differences of the quadratic and
quartic Casimirs. The system collapses to a $2\times2$ matrix built from one direct and one dual
Baxter polynomial on the two pole ladders of the measure. Its determinant vanishes for distinct
states and does not vanish for each state. The same holds with each row integrated over the
complete contour and a different pair of Q-functions in each row (Section~\ref{sec:full}). The
identities for the integrals of motion are proved
symbolically; the orthogonality and norm statements are verified exactly for every irrep at
$(S,P)=(1,1),(2,1),(2,2)$.
\end{abstract}

\section{Setup}

The chain has inhomogeneities $\theta_1,\theta_2$ and shift $h$, and $f\shift{\pm 2n}(u)=f(u\pm nh)$.
The Yangian weights on the highest-weight vector are
\begin{equation}
\nu_1=\nu_2=(u-\theta_1-Sh)(u-\theta_2-Ph),\qquad \nu_3=\nu_4=(u-\theta_1)(u-\theta_2).
\end{equation}
At zero twist ($z_j=1$) the transfer matrices $T_{a,1}=\tau_a$ have leading coefficient
$\binom{4}{a}$, and their subleading coefficient
\begin{equation}
[u^{2a-1}]\,\tau_a = -a\big(\theta_1+\theta_2+(a-1)h\big)\tbinom{4}{a} - h\tbinom{3}{a-1}\textstyle\sum_i\mu_i
\label{eq:subl}
\end{equation}
no longer depends on the individual weights $\mu_i$, only on $\sum_i\mu_i=2(S+P)$. As in the twisted
case, $\tau_3=\nu_1\shift{-4}\nu_4\,(4u^2+\cdots)$, $\tau_4$ is the quantum determinant, and five
coefficients remain free: $tc_1$ in $\tau_1$, $tc_{2,3,4}$ (the $u^2,u^1,u^0$ coefficients) in
$\tau_2$, and $tc_5$ in $\tau_3$.

The Baxter operator $\Op$ and its dual $\Op^\dagger$ are those of the twisted notes with $z_j=1$. The
measure is
\begin{equation}
\mu(u)=\prod_{\alpha=1}^2\prod_{n=0}^{S_\alpha-1}\frac{1}{u-\theta_\alpha-nh},\qquad S_1=S,\ S_2=P.
\end{equation}

\section{Q-functions at zero twist}

The irreps in $[S,S,0,0]\otimes[P,P,0,0]$ are $\mu=\{S+P-j,\,S+P-k,\,k,\,j\}$ with
$0\le j\le k\le\min(S,P)$, each with multiplicity one. For each irrep the direct solutions
$\Op\,q_k\shift{2}=0$ are monic polynomials of staggered degrees
\begin{equation}
\deg q_k = (S+P)-\mu_k+(k-1),\qquad k=1,\dots,4,
\end{equation}
and the dual solutions $\Op^\dagger q^j=0$ are monic of degrees
\begin{equation}
\deg q^4=\mu_4,\quad \deg q^3=\mu_3+1,\quad \deg q^2=\mu_2+2,\quad \deg q^1=\mu_1+3.
\end{equation}
In each family a lower solution may be added to a higher one. This freedom is removed by setting to
zero, in each polynomial, the coefficient of every lower degree of its family. With this gauge one
exact solve fixes the five $tc$ and the $q_k$ uniquely, and the $q^j$ then follow linearly and
uniquely. The Hodge-dual construction of the twisted case is not available, because its Vandermonde
normalisation vanishes at $z_j=1$. So the duals are obtained directly from $\Op^\dagger$.

\section{The orthogonality identity}

For eigenstates $A$ and $B$, $\Op^A$ annihilates $q^j_A$ from the left and $\Op^B$ annihilates
$q_i^{B[2]}$. With the contour enclosing all poles, the identity and quantum-determinant terms
cancel in the difference, which leaves
\begin{equation}
0=\oint du\ \mu\,q^j_A\Big[-\frac{\Delta T_1}{\nu_1}\,q^B_i
+\frac{\Delta T_2}{\nu_1\nu_1\shift{-2}}\,q_i^{B[-2]}
-\frac{\Delta T_3}{\nu_1\nu_1\shift{-2}\nu_1\shift{-4}}\,q_i^{B[-4]}\Big],
\qquad \Delta T_a=T^A_{a,1}-T^B_{a,1}.
\label{eq:identity}
\end{equation}
All poles lie on two ladders $\theta_\alpha+nh$, one per site. Inserting an $h$-periodic function
such as $e^{2\pi i u/h}$ leaves the derivation unchanged, and such a function is constant on each
ladder. So, for generic $\theta_1-\theta_2\notin h\mathbb Z$, \eqref{eq:identity} holds separately
for the sum of residues over each ladder. We write $L_s[\cdot]$, $s=1,2$, for the sum of residues
over the ladder of site $s$.

With twist, the right-hand side is linear in eight independent differences: the five $\Delta tc$
and three weight combinations entering the subleading coefficients. Distinct states make that
vector nonzero, and the $8\times8$ matrix of the twisted notes has vanishing determinant.

\section{The integrals of motion at zero twist}

At $z_j=1$ the eigenvalues depend on the irrep only through its Casimirs. With the shifted weights
$x_i=\mu_i+\tfrac52-i$ and $p_k=\sum_i x_i^k$,
\begin{equation}
C_2=p_2-5,\qquad
C_4=p_4+4p_3-p_1p_2+\tfrac12 p_2-2p_1^2-10p_1-\tfrac{51}{4}.
\end{equation}
The closed forms from \texttt{baxter\_SU4\_L2\_XXX.wb} are
\begin{align}
\tau_1&=\tau_1^{vac}+\tfrac{h^2}{2}\Delta C_2,\\
\tau_2&=\tau_2^{vac}+\tfrac{h^2}{8}\Delta C_2\big(\tau_1^{vac}+\tau_1^{vac[-2]}\big)
+h^4\big(\tfrac18\Delta(C_2^2)-\tfrac14\Delta C_4\big)+\tfrac{h^4(1+2S+2P)}{4}\Delta C_2,\\
\tau_3&=\nu_1\shift{-4}\nu_4\big[\,4u^2+\cdots+\tfrac{h^2}{2}\Delta C_2\,\big].
\end{align}
There $\Delta$ is measured from the vacuum $[S+P,S+P,0,0]$, $\Delta(C_2^2)$ denotes the difference
of the squares, and $\tau_1^{vac}=\sum_i\nu_i=2\nu_1+2\nu_4$. That notebook verified these forms
against explicitly constructed transfer matrices and against the Baxter solver. Here they reproduce
the Baxter-solved $tc$ of all twelve irreps at $(S,P)=(1,1),(2,1),(2,2)$.

For any two states it follows, identically in $u$ and symbolically in $\theta_1,\theta_2,h,S,P$
and two generic weights with $\sum\mu=2(S+P)$, that
\begin{equation}
\boxed{\;
\Delta T_1=\Delta tc_1,\qquad
\Delta T_2=\tfrac14\,\Delta tc_1\big(\tau_1^{vac}+\tau_1^{vac[-2]}\big)+\Delta c,\qquad
\Delta T_3=\nu_4\,\nu_1\shift{-4}\,\Delta tc_1\;}
\label{eq:ids}
\end{equation}
with the two independent invariants
\begin{equation}
\Delta tc_1=\tfrac{h^2}{2}\Delta C_2,\qquad
\Delta c=h^4\Big(\tfrac18\Delta(C_2^2)-\tfrac14\Delta C_4+\tfrac{1+2S+2P}{4}\Delta C_2\Big).
\label{eq:invariants}
\end{equation}
In terms of the five coefficients, \eqref{eq:ids} says
\begin{equation}
\Delta tc_5=\Delta tc_1,\qquad \Delta tc_2=2\,\Delta tc_1,\qquad
\Delta tc_3=-\big(2(\theta_1+\theta_2)+(S+P+2)h\big)\Delta tc_1,\qquad
\Delta tc_4=\Delta c+\tfrac14\Delta tc_1\big(\tau_1^{vac}(0)+\tau_1^{vac}(-h)\big).
\end{equation}
Together with \eqref{eq:subl}, which removes the three weight combinations, only two of the eight
differences of the twisted system stay independent. The dependent coefficients enter only through
the fixed combinations of $\nu_1$ and $\nu_4$ in \eqref{eq:ids}.

\section{The $2\times2$ matrix}

Substituting \eqref{eq:ids} into \eqref{eq:identity}, with
\begin{equation}
E_2=\frac{\mu}{\nu_1\nu_1\shift{-2}}\qquad\text{(poles at }\theta_\alpha+nh,\ n=0,\dots,S_\alpha+1),
\end{equation}
gives, for a fixed direct $q_i$ of $B$ and dual $q^j$ of $A$,
\begin{equation}
\boxed{\;
N^{AB}\cdot\begin{pmatrix}\Delta tc_1\\ \Delta c\end{pmatrix}=0,\qquad
\begin{aligned}
N^{AB}_{s,1}&=L_s\Big[E_2\,q^j_A\Big(-\nu_1\shift{-2}\,q^B_i
+\tfrac12\big(\nu_1+\nu_1\shift{-2}+\nu_4+\nu_4\shift{-2}\big)\,q_i^{B[-2]}
-\nu_4\,q_i^{B[-4]}\Big)\Big],\\
N^{AB}_{s,2}&=L_s\Big[E_2\,q^j_A\,q_i^{B[-2]}\Big],
\end{aligned}\;}
\label{eq:N}
\end{equation}
with $s=1,2$. Hence
\begin{equation}
\det N^{AB}=0\quad\text{whenever }(\Delta tc_1,\Delta c)\neq0,
\end{equation}
and $\det N^{AA}$ plays the role of the norm. No splitting of $T_{2,1}$ and no alignment of
measures is needed: each column is a single contour integral with the common factor $E_2$.

\section{Verification}

All statements below are exact. The identities \eqref{eq:ids}--\eqref{eq:invariants} are proved
symbolically. The Q-functions and matrices are computed with rational $\theta_1=\tfrac13$,
$\theta_2=\tfrac17$, $h=1$.

\paragraph{Q-functions.} For all twelve irreps of $(1,1)$, $(2,1)$, $(2,2)$ the gauge-fixed solve is
unique, $\Op\,q_k\shift{2}=0$ for $k=1,\dots,4$ and $\Op^\dagger q^j=0$ for $j=1,\dots,4$ hold
identically in $u$, and all irreps have distinct $tc$.

\paragraph{Choice of $(i,j)$.} Table~\ref{tab:survey} lists which pairs give a nonvanishing norm.
Eight pairs work for every irrep. The other eight give $\det N^{AA}=0$ for every irrep. Among the
failing pairs are the four pairings $i\in\{1,2\}$, $j\in\{3,4\}$ used in the twisted construction.

\begin{table}[h]
\centering
\begin{tabular}{ll}
\toprule
$\det N^{AA}\neq0$ for all 12 irreps & $(i,j)\in\{(2,1),(2,2),(3,1),(3,2),(3,3),(4,1),(4,2),(4,3)\}$\\
$\det N^{AA}=0$ for all 12 irreps & $(i,j)\in\{(1,1),(1,2),(1,3),(1,4),(2,3),(2,4),(3,4),(4,4)\}$\\
\bottomrule
\end{tabular}
\caption{Norm survey over the direct index $i$ and dual index $j$.}
\label{tab:survey}
\end{table}

\paragraph{Orthogonality and norms.} At $(i,j)=(2,1)$, for every ordered pair of irreps:

\begin{center}
\begin{tabular}{lc}
\toprule
norms with $\det N^{AA}\neq0$ & 12 / 12\\
cross pairs with $\det N^{AB}=0$ & 42 / 42\\
cross pairs with $N^{AB}(\Delta tc_1,\Delta c)^T=0$ & 42 / 42\\
\quad of which $N^{AB}\equiv0$ & 6\\
\bottomrule
\end{tabular}
\end{center}

\paragraph{Example, $(S,P)=(1,1)$, $(i,j)=(2,1)$.} The polynomials entering are
$q_2^{\{2,1,1,0\}}=u^2-\tfrac{10}{21}u$,
$q^1_{\{2,1,1,0\}}=u^5-\tfrac{130}{21}u^4+\tfrac{3097}{63}u$ and
$q^1_{\{2,2,0,0\}}=u^5-\tfrac{22525}{1323}u^3+\tfrac{16640}{441}u^2$. Then
\begin{equation}
N^{AA}\big|_{A=\{2,1,1,0\}}=
\begin{pmatrix}-\frac{150144211}{334305} & \frac{3235736}{334305}\\[3pt]
\frac{692896376}{1820105} & -\frac{41930251}{1820105}\end{pmatrix},\qquad
\det N^{AA}=\frac{1039313041}{156009},
\end{equation}
\begin{equation}
N^{AB}\big|_{A=\{2,2,0,0\},\,B=\{2,1,1,0\}}=
\begin{pmatrix}-\frac{320542493}{5014575} & \frac{320542493}{5014575}\\[3pt]
\frac{1737705523}{27301575} & -\frac{1737705523}{27301575}\end{pmatrix},\qquad
\det N^{AB}=0,
\end{equation}
and $(\Delta tc_1,\Delta c)=(2,2)$ spans the kernel of $N^{AB}$.

\section{Variant: full-contour rows with different Q-functions}
\label{sec:full}

Instead of splitting the contour by ladder, one may integrate over the complete contour (both
ladders, no periodic insertion) and use a different pair of Q-functions $(i_r,j_r)$ in each row
$r=1,2$. With $\oint=L_1+L_2$,
\begin{equation}
\tilde N^{AB}_{r,1}=\oint E_2\,q^{j_r}_A\Big(-\nu_1\shift{-2}\,q^B_{i_r}
+\tfrac12\big(\nu_1+\nu_1\shift{-2}+\nu_4+\nu_4\shift{-2}\big)\,q_{i_r}^{B[-2]}
-\nu_4\,q_{i_r}^{B[-4]}\Big),\qquad
\tilde N^{AB}_{r,2}=\oint E_2\,q^{j_r}_A\,q_{i_r}^{B[-2]},
\label{eq:Nfull}
\end{equation}
and again $\tilde N^{AB}\cdot(\Delta tc_1,\Delta c)^T=0$. Each row is the sum of the two ladder rows
of \eqref{eq:N}. Because the contour encloses every finite pole, a row equals minus the residue at
infinity of its integrand, so it can vanish identically when the degrees of the Q-functions are low.

Exactly, with the same parameters and irreps:
\begin{itemize}
\item A full-contour row vanishes identically for every irrep exactly when $i<j$. This includes
the four twisted pairings $i\in\{1,2\}$, $j\in\{3,4\}$.
\item Among the ten remaining rows, $(1,1)$ and $(4,4)$ give a vanishing norm when paired, and so
does any two of $(2,2),(3,2),(3,3)$. For the norm these rows are therefore linearly dependent.
\item Every other pair of nonvanishing rows gives $\det\tilde N^{AA}\neq0$ for all twelve irreps.
This is 41 of the 120 pairs of distinct rows. The other 79 pairs give $\det\tilde N^{AA}=0$ for all
twelve irreps; no pair works for only some irreps.
\end{itemize}
At the rows $\big((i_1,j_1),(i_2,j_2)\big)=\big((1,1),(2,1)\big)$, that is $q^1_A$ against $q^B_1$
and against $q^B_2$:

\begin{center}
\begin{tabular}{lc}
\toprule
norms with $\det\tilde N^{AA}\neq0$ & 12 / 12\\
cross pairs with $\det\tilde N^{AB}=0$ & 42 / 42\\
cross pairs with $\tilde N^{AB}(\Delta tc_1,\Delta c)^T=0$ & 42 / 42\\
\quad of which $\tilde N^{AB}\equiv0$ & 6\\
\bottomrule
\end{tabular}
\end{center}

The row $(1,1)$ is usable here, although on a single ladder $(1,1)$ gives a vanishing norm
(Table~\ref{tab:survey}).

\section{Remarks and open points}

\begin{enumerate}
\item \textbf{Vacuous pairs.} For six ordered cross pairs, all at $(S,P)=(2,2)$, the matrix
vanishes identically. They form three unordered pairs:
$\{4,4,0,0\}\leftrightarrow\{2,2,2,2\}$, $\{4,3,1,0\}\leftrightarrow\{3,2,2,1\}$ and
$\{4,2,2,0\}\leftrightarrow\{3,3,1,1\}$. For these pairs \emph{every} row vanishes: all 16
full-contour rows and all 32 single-ladder rows, for every $(i,j)$. So neither \eqref{eq:N} nor
\eqref{eq:Nfull} tests orthogonality there, for any choice of Q-functions or contour. The vanishing
is structural, like a selection rule, rather than an artefact of the choice. In the notation
$\{4-j,4-k,k,j\}$ each pair satisfies $(j+k)_A+(j+k)_B=4$. The reason is not understood, and it is
not a Casimir degeneracy ($\Delta C_2\neq0$).
\item \textbf{Working $(i,j)$.} The pairs with nonvanishing norm are those with $i\ge j$, except
$(1,1)$ and $(4,4)$. This is observed at three spins and not derived.
\item \textbf{Degenerate states.} At zero twist the chain is $\mathfrak{gl}(4)$ symmetric.
Descendants share the transfer-matrix eigenvalues of their highest weight, so $\Delta tc_1=\Delta c=0$
and the argument separates only one state per irrep.
\item \textbf{Relation to the twisted construction.} The three weight differences that disappear
at $z_j=1$ are exactly the Cartan charges $\sum_i\mu_i z_i\partial_{z_i}\chi_a$ carried by the
subleading coefficients with twist. The twisted pairings $i\in\{1,2\}$, $j\in\{3,4\}$ each give
vanishing $2\times2$ norms at zero twist.
\end{enumerate}

\end{document}
